% Generated draft: requires mathematical and source review.
\documentclass[11pt]{article}
\usepackage[margin=25mm]{geometry}
\usepackage{fontspec}
\setmainfont{Latin Modern Roman}
\usepackage{amsmath,amssymb,mathtools,unicode-math}
\setmathfont{Latin Modern Math}
\usepackage{enumitem,needspace,xurl,hyperref}
\setlength{\parindent}{0pt}
\setlength{\parskip}{5pt}
\setlength{\emergencystretch}{3em}
\newcommand{\N}{\mathbb{N}}
\newcommand{\Z}{\mathbb{Z}}
\newcommand{\Q}{\mathbb{Q}}
\newcommand{\R}{\mathbb{R}}
\newcommand{\C}{\mathbb{C}}
\newcommand{\F}{\mathbb{F}}
\newcommand{\A}{\mathbb{A}}
\newcommand{\PP}{\mathbb{P}}
\newcommand{\E}{\mathbb{E}}
\newcommand{\eps}{\varepsilon}
\newcommand{\dd}{\,\mathrm{d}}
\begin{document}
\section*{UnsolvedMath ID: 10900030}
\textbf{Source problem code:} catalogue entry
\subsection{Definitions and mathematical statement}
Fix $0<\mu<\mu_3$, where $\mu_3$ is a three-dimensional Margulis constant. For a complete finite-volume hyperbolic $3$-manifold $M$ with $\operatorname{Vol}(M)<V$, let $D_\mu(M)$ be the manifold obtained by drilling all closed geodesics of length strictly less than $\mu$. Drilling means removing tubular neighborhoods of these geodesics, or equivalently considering the corresponding cusped interior. Let
\[
 \mathcal D_\mu(V)=\{D_\mu(M):\operatorname{Vol}(M)<V\}/\cong
\]
be the distinct resulting manifolds, with duplicate outputs counted once. The source states that this collection is finite. The volume restriction applies to $M$ before drilling, not to the output.

A resulting manifold is small if it contains no closed incompressible non-boundary-parallel surface, in the usual essential-surface convention. For nonempty $\mathcal D_\mu(V)$ define
\[
 s_\mu(V)=
 \frac{\#\{N\in\mathcal D_\mu(V):N\text{ is small}\}}
 {\#\mathcal D_\mu(V)}.
\]
Agol's question, item 5.5, asks for the limiting behavior of $s_\mu(V)$ as $V\to\infty$ and its dependence on $\mu$. It does not assert that a limit exists: determining existence or alternative asymptotic behavior is part of the request.

The source does not state whether the input class is closed only or also includes cusped manifolds, nor its orientability and duplicate-identification conventions. The displayed complete finite-volume model records the usual setting without resolving those missing choices. Smallness concerns embedded closed essential surfaces and must not be replaced by absence of immersed $\pi_1$-injective surfaces.
\subsection{Short English statement}
After drilling all geodesics shorter than a fixed Margulis threshold and counting distinct outputs once, how does the fraction of small manifolds behave as the original volume cutoff grows, and how does the threshold affect it?
\subsection{Sources}
\begin{itemize}
\item Local registry: \url{data/unsolvedmath/problems.json}, permanent UnsolvedMath ID $10900030$, exact source code \texttt{catalogue entry}, attributed to Agol.
\item Delp, Hoffoss and Manning, \emph{Problems In Groups, Geometry, and Three-Manifolds} (2015), subsection 5.5, PDF page 4: \url{https://arxiv.org/pdf/1512.04620#page=4}; \url{https://arxiv.org/abs/1512.04620}. Supplied access date: 2026-07-29; source-tex extraction checked against rendered PDF text by the source record.
\item The registry mentions related asymptotic volume/geodesic counts without precise references, and reports no verified determination of the limiting fraction.
\end{itemize}
\end{document}
